\chapter{Decidir bajo incertidumbre}\label{ch:decisions-under-uncertainty}

% Alcance: valor esperado, Bayes, árboles de decisión, sesgos cognitivos.
% Vínculos cruzados: la actualización bayesiana reaparece en atribución basada en datos y MMM (\cref{ch:attribution-and-measurement});
%   la regla de la cadena de probabilidad subyace al funnel de conversión (\cref{ch:digital-funnel-and-crm}).

Toda decisión de negocio es una apuesta colocada antes de conocer el resultado. Los \emph{framework}s
aquí presentados --- valor esperado, Bayes y teoría prospectiva --- no eliminan la incertidumbre; hacen
explícita la estructura de la apuesta para que el directivo pueda evaluarla con claridad en lugar de
reaccionar emocionalmente ante ella.

\section{Valor esperado y árboles de decisión}\label{sec:expected-value}
% complejidad: 4 -> figura decision-tree

Cuando dos opciones tienen perfiles de riesgo distintos, ¿cuál se elige? El valor esperado convierte
los resultados ponderados por probabilidad en un único número comparable, pero un árbol de decisión
obliga a hacer explícita la estructura de la incertidumbre antes de calcular. Comparar una opción
segura con una arriesgada usando el juicio intuitivo induce sistemáticamente al error, porque los
seres humanos sobrepondera los resultados recientes, vívidos y emocionalmente salientes. El valor
esperado es el promedio aritmético de largo plazo --- el número hacia el que se convergería si la
decisión se tomara muchas veces en condiciones idénticas. Para decisiones únicas, EV es una línea
base disciplinada, no una garantía.

Para un conjunto discreto de resultados mutuamente excluyentes \(x_i\) con probabilidades \(p_i\)
(\(\sum_i p_i = 1\)):
\begin{equation}\label{eq:ev}
  \mathbb{E}[X] \;=\; \sum_i p_i\, x_i.
\end{equation}
Un árbol de decisión mapea cada elección a un conjunto de ramas ponderadas por probabilidad, con EV
calculado en cada nodo terminal y enrollado de vuelta al nodo de decisión mediante \cref{eq:ev}. El
valor del árbol radica en obligar al analista a enumerar los resultados y asignar probabilidades, en
lugar de razonar a partir de un único escenario esperado.

\begin{figure}[htbp]
  \centering
  \includegraphics[width=0.86\linewidth]{decision-tree}
  \caption{Árbol de decisión para dos pujas de CPA objetivo. La puja más alta (\money{900}) tiene mayor
    valor esperado pese a una probabilidad de éxito menor, porque el piso de pérdida es pequeño en
    relación con el techo de ganancia. Enrollar el árbol hacia el nodo de decisión revela esto sólo
    después de hacer explícitas las probabilidades.}
  \label{fig:decision-tree}
\end{figure}

\Cref{eq:ev} supone que los resultados y las probabilidades son conocidos. En la práctica, ambos
requieren estimación a partir de datos o juicio experto, lo que introduce riesgo de modelo. EV
tampoco capta la distribución alrededor de la media: dos opciones con EV idéntico pueden tener
varianza muy diferente, y una empresa con capital restringido puede preferir racionalmente la opción
de menor varianza aunque su EV sea inferior (aversión al riesgo --- formalizada en
\cref{sec:prospect-theory}). Por último, el EV de un árbol de decisión colapsa a un único número y
descarta la temporalidad de los \emph{cash flow}s; se debe incorporar descuento al comparar trayectorias
multiperiódicas.

\begin{example}
La empresa SaaS de referencia evalúa dos niveles de CPA objetivo para una nueva campaña. Con
\(\mathrm{CPA}=\money{600}\), el modelo de conversión asigna una probabilidad del \qty{70}{\percent} de
adquirir un cliente con \(\mathrm{LTV}=\money{3150}\) y \qty{30}{\percent} de un cliente de bajo
compromiso valorado en \money{900} (tres meses y luego \emph{churn}):
\[
  \begin{aligned}
    \mathbb{E}[\text{value} \mid \mathrm{CPA}=\money{600}]
      &= 0.70\times(\money{3150}-\money{600}) + 0.30\times(\money{900}-\money{600}) \\
      &= \money{1785} + \money{90} = \money{1875}.
  \end{aligned}
\]
Con \(\mathrm{CPA}=\money{900}\), el mayor gasto atrae un segmento de mayor intención:
\qty{85}{\percent} de probabilidad del cliente de LTV completo, \qty{15}{\percent} del de bajo compromiso:
\[
  \begin{aligned}
    \mathbb{E}[\text{value} \mid \mathrm{CPA}=\money{900}]
      &= 0.85\times(\money{3150}-\money{900}) + 0.15\times(\money{900}-\money{900}) \\
      &= \money{1912.50}.
  \end{aligned}
\]
El CPA de \money{900} genera mayor beneficio esperado por cliente (\money{1912.50} frente a
\money{1875}) pese a gastar \money{300} más por adquisición. La decisión: subir el CPA objetivo. La
trampa es elegir \money{600} porque \emph{parece} más conservador --- \cref{sec:prospect-theory}
explica por qué esa intuición es sistemáticamente errónea.
\end{example}

\begin{admonition}{Advertencia}
La planificación de un único escenario --- construir un conjunto de proyecciones, denominarlo
«caso base» y evaluarlo como si fuera cierto --- es un árbol de decisión con una sola rama, que
asigna probabilidad~1 al escenario elegido. Hacer explícitas las ramas y sus probabilidades no es
trabajo adicional; es la disciplina mínima que exige el análisis de EV \parencite{kahneman2011}.
\end{admonition}

Cuando las decisiones son secuenciales y la información llega entre etapas, el árbol se convierte en
un problema de opciones reales: el valor de poder esperar, expandir o abandonar es una prima de opción
sobre el EV estático. La valoración de Black-Scholes es el límite en tiempo continuo de un árbol
binomial. La implicación práctica: escalonar los compromisos de capital (primero un piloto, luego el
despliegue completo) casi siempre vale más que su costo, porque preserva la opción de abandonar.
Modelar esto requiere métodos de valoración de opciones, no EV simple (\cref{ch:valuation}).

El valor esperado es el lenguaje de la presupuestación y el análisis de escenarios; alimenta directamente
el \emph{framework} DCF de \cref{ch:valuation} y los modelos de atribución MMM de
\cref{ch:attribution-and-measurement}. \textcite{anderson2019statistics} cubren el aparato probabilístico;
\textcite{kahneman2011} documentan por qué los directivos se desvían de él.

\section{Actualización bayesiana}\label{sec:bayesian-updating}
% complejidad: 3

Ha llegado una nueva señal --- la descarga de un documento técnico, un resultado de piloto, una señal
de \emph{churn}. ¿Cuánto debe mover la creencia previa? El teorema de Bayes da la respuesta exacta;
la habilidad práctica está en elegir qué cuenta como evidencia y calibrar su poder diagnóstico.
La inferencia frecuentista pregunta cuál es la probabilidad de estos datos dado un parámetro verdadero
fijo; la inferencia bayesiana pregunta cuál es la probabilidad de que el parámetro sea verdadero dados
los datos recién observados. La segunda pregunta es casi siempre la que necesita el directivo. Bayes
dice: comience con una prior (la creencia inicial), observe la evidencia y actualice proporcionalmente
a cuánto más probable es la evidencia bajo la hipótesis que bajo sus alternativas.

Sea \(H\) una hipótesis y \(D\) los datos observados. La probabilidad posterior de \(H\) dado \(D\) es:
\begin{equation}\label{eq:bayes}
  P(H \mid D) \;=\; \frac{P(D \mid H)\,P(H)}{P(D)},
\end{equation}
donde \(P(H)\) es la prior, \(P(D \mid H)\) la verosimilitud de los datos bajo \(H\), y
\(P(D)=P(D\mid H)P(H)+P(D\mid \neg H)P(\neg H)\) la constante normalizadora (verosimilitud marginal).
La actualización es multiplicativa en log-odds: \(\log\frac{P(H\mid D)}{P(\neg H\mid D)} =
\log\frac{P(H)}{P(\neg H)} + \log\frac{P(D\mid H)}{P(D\mid \neg H)}\). Cada nueva observación
agrega su razón de log-verosimilitud al log-odds acumulado.

\Cref{eq:bayes} es exacta, pero sólo tan buena como sus insumos. Una prior mal calibrada (p.\ ej.,
una tasa base ignorada en favor del caso vívido) se propaga a la posterior. Una verosimilitud
construida sobre datos históricos no representativos distorsiona igualmente la actualización. La
fórmula no puede corregir insumos deficientes; sólo puede hacer transparente la estructura del
razonamiento para que los insumos puedan cuestionarse.

\begin{example}
El equipo de ventas estima que el \qty{25}{\percent} de los usuarios en periodo de prueba acaba
convirtiéndose a pago --- la prior: \(P(\text{paid})\!=\!0.25\). Se observa la descarga de un
documento técnico. Los datos históricos muestran que el \qty{60}{\percent} de los usuarios que
eventualmente pagaron descargaron un documento técnico durante el periodo de prueba, frente a sólo
el \qty{15}{\percent} de los usuarios que hicieron \emph{churn}: \(P(D\mid\text{paid})=0.60\),
\(P(D\mid\text{not paid})=0.15\). Aplicando \cref{eq:bayes}:
\[
  P(D) = 0.60\times0.25 + 0.15\times0.75 = 0.15 + 0.1125 = 0.2625,
\]
\[
  P(\text{paid}\mid D) = \frac{0.60\times0.25}{0.2625} \approx 0.571.
\]
Una única descarga de documento técnico más que duplica la probabilidad de conversión estimada, del
\qty{25}{\percent} al \qty{57}{\percent}. La implicación para el CRM: dirigir este \emph{lead} a una
secuencia de ventas de mayor dedicación cuyo costo queda justificado por la estimación de LTV
revisada. Sin la actualización bayesiana, ambos \emph{lead}s reciben un tratamiento idéntico pese a
tener valores esperados muy distintos.
\end{example}

\begin{admonition}{Advertencia}
La negligencia de la tasa base es el error bayesiano más común: actualizar con la señal vívida
ignorando la prior. Una \emph{startup} anuncia que su tasa de conversión de prueba a pago es del
\qty{40}{\percent} --- ¿pero eso es condicional a tener el mismo \emph{funnel} de adquisición y la
misma calidad de producto que la referencia? Sin una prior calibrada y una razón de verosimilitud
cuidadosamente medida, el \qty{40}{\percent} contamina la posterior en lugar de informarla
\parencite{kahneman2011}.
\end{admonition}

Las pruebas A/B bayesianas reemplazan el \emph{framework} frecuentista de hipótesis nula por una
posterior sobre la diferencia en la tasa de conversión. Esto genera una cantidad relevante para la
decisión («¿cuál es la probabilidad de que la variante B sea mejor?») en lugar de un valor p («¿qué
tan probables son estos datos si la hipótesis nula es verdadera?»). Los algoritmos de bandido
multi-brazo extienden el \emph{framework} para reasignar continuamente el tráfico hacia las variantes
de mejor rendimiento, haciendo explícito el intercambio exploración--explotación. Los modelos de
atribución de \cref{ch:attribution-and-measurement} son bayesianos en su núcleo: actualizan los
pesos de atribución de canal a medida que llegan nuevos datos de conversión.

La actualización bayesiana es la columna vertebral de la atribución probabilística en
\cref{ch:attribution-and-measurement} y de los modelos de puntuación de clientes que alimentan el
\emph{funnel} en \cref{ch:digital-funnel-and-crm}. \textcite{anderson2019statistics} cubren los
\emph{framework}s frecuentista y bayesiano en paralelo.

\section{Sesgo cognitivo y teoría prospectiva}\label{sec:prospect-theory}
% complejidad: 4 -> figura prospect-value-fn

El equipo tiene dos opciones con valor esperado idéntico. Elegirá de manera fiable la menos
arriesgada. La teoría prospectiva explica las formas específicas y predecibles en que la percepción
del riesgo humana se desvía de la maximización de EV, para poder corregirlo. Las personas evalúan
los resultados en relación con un punto de referencia (habitualmente el statu quo), no en términos
absolutos. Las ganancias se perciben positivamente pero con sensibilidad decreciente; las pérdidas
desde ese mismo punto de referencia se perciben como aproximadamente el doble de malas que las
ganancias equivalentes se perciben de buenas. La función de valor resultante tiene forma de S:
cóncava en el dominio de ganancias (aversión al riesgo), abruptamente convexa en el dominio de
pérdidas (búsqueda de riesgo para evitar una pérdida segura). Esto explica por qué un CPA seguro de
\money{600} «parece más seguro» que una puja de \money{900} con mayor beneficio esperado: el gasto
adicional de \money{300} se codifica como una pérdida segura, ponderada a aproximadamente el doble
de su valor nominal.

\textcite{kahneman1979prospect} especifican la función de valor de la teoría prospectiva como:
\begin{equation}\label{eq:prospect}
  v(x) \;=\;
  \begin{cases}
    x^{\alpha} & x \ge 0, \\[3pt]
    -\lambda\,(-x)^{\beta} & x < 0,
  \end{cases}
\end{equation}
donde \(x\) es el resultado relativo al punto de referencia, \(\alpha,\beta\in(0,1)\) capturan la
sensibilidad decreciente en ambos dominios (\(\alpha\approx\beta\approx0.88\) empíricamente), y
\(\lambda>1\) es el \emph{coeficiente de aversión a la pérdida} (\(\lambda\approx2.25\) en promedio).
La función no es una función de utilidad sobre la riqueza; es una función de valor sobre cambios.
La teoría prospectiva también reemplaza las probabilidades objetivas por pesos de decisión que
sobrepondera las probabilidades pequeñas y subpondera las moderadas a grandes, pero la función de
valor impulsa la mayoría de las predicciones relevantes para la gestión.

\begin{figure}[htbp]
  \centering
  \includegraphics[width=0.86\linewidth]{prospect-value-fn}
  \caption{La función de valor de la teoría prospectiva (\cref{eq:prospect}, \(\lambda=2.25\)). Cóncava
    en ganancias (aversión al riesgo); abruptamente convexa en pérdidas (búsqueda de riesgo para
    evitar una pérdida segura). La asimetría en cero significa que una pérdida de cualquier
    magnitud pesa más que una ganancia equivalente.}
  \label{fig:prospect-value-fn}
\end{figure}

\Cref{eq:prospect} es un modelo descriptivo del comportamiento humano promedio en condiciones de
laboratorio, no una regla normativa. No se sostiene de manera uniforme entre culturas, niveles de
riesgo o dominios con experiencia de juego repetido. El punto de referencia en sí mismo se desplaza:
una bonificación enmarcada como evitar una penalización genera elecciones diferentes que la misma
bonificación enmarcada como una ganancia, incluso cuando el resultado objetivo es idéntico (el
\emph{efecto de encuadre}). Los directivos que comprenden esto pueden diseñar estructuras de
compensación e informes de experimentos que reduzcan el sesgo, pero no deben asumir que
\(\lambda=2.25\) es una constante universal.

\begin{example}
Retomando las dos pujas de CPA. El CPA de \money{600} parece más seguro porque pagar \money{900}
se codifica como una pérdida adicional de \money{300} respecto al punto de referencia de \money{600}.
Aplicando \cref{eq:prospect} con \(\lambda=2.25\):
\[
  v(-\money{300}) \approx -2.25\times(300)^{0.88} \approx -2.25\times191 \approx -430
  \quad\text{(pérdida)}
\]
frente a la ganancia del mayor rendimiento esperado:
\[
  v(+\money{37.50}) \approx (37.50)^{0.88} \approx 27
  \quad\text{(ganancia por diferencia de EV \money{1912.50} vs.\ \money{1875})}.
\]
El costo psíquico de la puja mayor (\(-430\) unidades de valor) supera ampliamente la ganancia
psíquica del EV más alto (+27 unidades) --- de modo que el directivo siente racionalmente que la
puja de \money{600} es mejor, aunque el beneficio esperado sea mayor a \money{900}. La corrección:
anclar la decisión al EV (\cref{eq:ev}), no a la magnitud del gasto incremental.
\end{example}

\begin{admonition}{Advertencia}
La falacia de la planificación --- subestimar el costo y la duración de un proyecto porque el punto
de referencia es el plan, no las tasas base históricas de proyectos similares --- amplifica la
aversión a la pérdida: los directivos son simultáneamente demasiado confiados sobre los resultados
y excesivamente reacios a los escenarios negativos que no pueden imaginar vívidamente
\parencite{tversky1974judgment}. Los premortem, en los que el equipo imagina que el proyecto ya
fracasó y pregunta por qué, son la técnica de desesgado más confiablemente eficaz.
\end{admonition}

El estatus normativo de la teoría prospectiva es objeto de debate. \textcite{kahneman1979prospect}
la presentan como un modelo descriptivo; los teóricos de la utilidad esperada (\cref{eq:ev}) argumentan
que la desviación de la maximización de EV es simplemente un error de decisión que mejores incentivos
o educación pueden eliminar. Los economistas del comportamiento responden que los sesgos son demasiado
estables y demasiado grandes para ser mero ruido. La resolución práctica: usar EV como criterio de
decisión, pero diseñar la arquitectura de la elección --- cómo se enmarcan, secuencian y presentan
las opciones --- para reducir la brecha entre las evaluaciones de EV y las de la teoría prospectiva
\parencite{kahneman2011}.

\textcite{thaler1985mental} introduce la contabilidad mental --- la tendencia a evaluar los resultados
en cuentas psicológicas separadas en lugar de como cambios en la riqueza total. Una implicación para
el \emph{marketing}: incluir un aumento de precio junto con una mejora visible del producto traslada
el cambio neto de la cuenta de pérdidas a la cuenta de «intercambio justo», reduciendo la pérdida
percibida. Los efectos de encuadre y la contabilidad mental son el mecanismo que subyace al anclaje
(ejemplo resuelto de \cref{sec:prospect-theory}) y a las estrategias de partición de precios
exploradas en \cref{ch:pricing}.

La teoría prospectiva es el fundamento de la arquitectura de precios en \cref{ch:pricing} y de la
comprensión de por qué los clientes resisten el cambio incluso cuando la alternativa tiene mayor
valor esperado. Los \emph{framework}s bayesiano y prospectivo sustentan conjuntamente la discusión
sobre atribución en \cref{ch:attribution-and-measurement}. \textcite{tversky1974judgment} documentan
las heurísticas; \textcite{kahneman1979prospect} formalizan la función de valor;
\textcite{kahneman2011} sintetizan ambas en un \emph{framework} práctico.
